\section{Discussion}

This section summarizes the scope and limits of the framework. The aim is not to add new claims, but to clarify what the framework does and does not address.

\subsection{Scope of Description}

The paper describes biological systems using information density \(\I\), information flow \(\J\), and induced potential \(\PhiI\). This is a minimal language for treating development, morphogenesis, life maintenance, and the conditions for intelligence in a common form.

The framework extracts structure. It does not attempt detailed reproduction of each individual biological phenomenon.

\subsection{Relation to Existing Models}

The framework does not replace existing biological models. Reaction--diffusion systems, gene networks, and neural dynamics remain necessary at their own levels. The present formalism provides a way to redescribe such models in terms of common structural elements: density, flow, induced bias, and stability.

It does not require the assumption of a new mechanism, nor does it deny existing theories.

\subsection{Non-Closure and Sustained Dynamics}

The dynamics considered here do not generally converge to complete equilibrium or fixed closure. Flow continues, and structure is updated. In this sense, biological systems are interpreted as sustained non-closed structures.

This property is not developed as a separate theory in this paper. It appears as a descriptive feature of the dynamics considered.

This interpretation is consistent with the quasi-closure structure formalized in IFGT: biological systems remain within regimes where closure is partial, dynamically maintained, and not realized as an ordinary complete-closure state. The non-closure of biological dynamics is therefore not a limitation of the framework, but its central descriptive feature.

\subsection{Status of Time}

The paper uses external time \(t\) and introduces internal time \(\tau\) as an interpretive variable. This separates observed time from the progress of internal state update. A precise operational definition of \(\tau\) is outside the scope of the present paper.

\subsection{Limits}

Several limits should be stated. The framework does not derive the parameter set \(\theta\) for a concrete organism. The explicit forms of \(\I\), \(\J\), and \(\PhiI\) are system-dependent. The paper does not provide numerical predictions.

The framework is therefore a unifying description, not a detailed model of a particular biological system.

\subsection{Future Directions}

Future work should identify \(\I\), \(\J\), and \(\PhiI\) in concrete systems, analyze the structure of parameter spaces, and quantify attractor regions. For intelligence, further work must address memory, temporal integration, language, and interaction among individuals.

\subsection{Summary}

The framework describes biological phenomena in a single formal language without using design or purpose as primitives. It does not address all mechanistic detail or provide numerical prediction. It is best understood as a basis for description and later specialization.
